\chapter{Oh My Pi}
\chapterlead{Oh My Pi（命令名 \texttt{omp}）是一个终端内的编码 Agent，提供语言服务、调试器接入、会话树与多模型切换。本章说明安装、基本操作与 NixOS 上的声明式配置。}

\section{版本与前提}

截至 2026 年 11 月 25 日：

\begin{booktable}{\textwidth}{J{0.86} J{1.14}}
\toprule
项目 & 值 \\
\midrule
npm 包 & \texttt{@oh-my-pi/\allowbreak{}pi-coding-agent} \\
版本 & 19.1.1 \\
运行时 & Bun 1.4.4 或更新 \\
命令 & \texttt{omp} \\
安装脚本 & \texttt{omp.sh/\allowbreak{}install} \\
\bottomrule
\end{booktable}

\section{安装}

官方提供四条路径。NixOS 上用第三或第四条，避免把运行时装进系统。

\begin{shellsession}
# 一、官方脚本，下载预编译二进制并校验 SHA-256
$ curl -fsSL https://omp.sh/install | sh

# 二、Bun
$ bun install -g @oh-my-pi/pi-coding-agent

# 三、Nix，直接运行
$ nix run github:can1357/oh-my-pi
\end{shellsession}

在 NixOS 配置中以模块方式安装，配置也由 Nix 管理：

\begin{nixcode}
{
  inputs.omp.url = "github:can1357/oh-my-pi";

  # 在 Home Manager 模块中
  imports = [ inputs.omp.homeManagerModules.default ];
  programs.omp = {
    enable = true;
    settings.startup.quiet = true;
  };
}
\end{nixcode}

这份配置把 \texttt{omp} 装进用户环境，并把设置写进 \texttt{settings.json}，随配置重建生效。

\section{第一次会话}

\begin{shellsession}
$ cd ~/proj
$ omp
$ omp --continue        # 继续上一次会话
$ omp --resume          # 从历史会话中选择
\end{shellsession}

进入后直接描述目标。界面底部显示当前模型与快捷键提示。会话自动保存在 \texttt{\textasciitilde/.omp/agent/sessions} 下，按工作目录分组，格式为 JSONL，一行一条消息。

常用操作：

\begin{booktable}{\textwidth}{J{0.692} J{1.308}}
\toprule
操作 & 用途 \\
\midrule
\texttt{Ctrl+L} & 切换模型 \\
\texttt{/compact} & 压缩上下文 \\
\texttt{/tree} & 浏览会话分支 \\
\texttt{!命令} & 在会话内执行命令并把输出加入上下文 \\
\texttt{Esc} & 中断当前生成 \\
\bottomrule
\end{booktable}

\texttt{!} 前缀值得单独说明：输入 \texttt{!git status} 会把结果直接并入上下文，比自己跑完再粘贴更省事，也更准确。

\section{AGENTS.md}

Agent 启动时读取项目说明文件，顺序是：全局 \texttt{\textasciitilde/.omp/agent/AGENTS.md}、从当前目录向上各级目录、当前目录。写清构建命令、测试命令、代码风格与禁止事项，能显著减少来回确认。

\begin{textcode}
# 项目说明

## 常用命令
- 构建：pnpm build
- 测试：pnpm test
- 检查：pnpm lint

## 约定
- TypeScript 严格模式
- 提交信息用中文，标题不超过 50 字
- 不要改动 migrations 目录
\end{textcode}

这份文件是给 Agent 的，也是给新同事的。能被写下来的约定，比口头传达可靠。

\section{语言服务与调试}

\texttt{omp} 接入了语言服务与调试器协议，因此重命名、跳转、查看类型这类操作走的是编译器而不是文本替换。前提是项目里对应的语言服务可用。

\begin{shellsession}
$ omp --help | grep -i lsp
$ which typescript-language-server
\end{shellsession}

语言服务缺失时，Agent 会退回文本搜索，改名的可靠性下降。为项目准备环境时，把语言服务一并装进 devShell，Agent 和编辑器都能受益。

\section{成本与模型}

不同模型的单价差别很大。两件事能显著降低成本：把简单任务交给小模型，把上下文控制在必要范围。

\begin{shellsession}
$ omp --model <模型名>
$ omp --smol <小模型名>     # 简单任务使用
$ omp --plan                # 只规划，不执行
\end{shellsession}

长会话的上下文会不断增长，每次请求都携带全部历史，费用随之上升。一个任务完成后开新会话，比在一个会话里做十件事便宜。

\section{一次完整任务}

目标：给一个函数补单元测试，并让覆盖率上升到 80\% 以上。

\begin{enumerate}
  \item 确认 \texttt{AGENTS.md} 里写明了测试命令。
  \item 让它先读实现与现有测试，输出计划，不改文件。
  \item 确认计划后放开执行，观察它是否只改测试文件。
  \item 跑测试与覆盖率，失败时把输出贴回会话。
  \item 用 \texttt{git diff} 检查是否动到了无关文件。
\end{enumerate}

\begin{shellsession}
$ git diff --stat
$ pnpm test -- --coverage
\end{shellsession}

\section{练习}

\begin{enumerate}
  \item 用 \texttt{nix run} 运行一次，记录版本号。
  \item 在项目里写一份 \texttt{AGENTS.md}，包含三条命令与两条约定。
  \item 用 \texttt{!} 前缀执行一条命令，确认输出进入了上下文。
  \item 用 \texttt{/tree} 回到会话中间某一轮，从那里继续，说明分支的行为。
  \item 用 \texttt{-{}-plan} 跑一个任务，比较它与直接执行在成本与结果上的差别。
\end{enumerate}

\begin{checkpoint}
本章结束时，应能：
\begin{itemize}
  \item 说出当前版本、运行时要求与四条安装路径；
  \item 用 Nix 模块声明式安装并管理设置；
  \item 用 \texttt{--continue}、\texttt{--resume}、\texttt{/tree} 管理会话；
  \item 写出有用的 \texttt{AGENTS.md}；
  \item 通过模型选择与会话拆分控制成本。
\end{itemize}
\end{checkpoint}
